\section{What matrices are}

A matrix is a rectangular arrangement of mathematical objects, usually numbers. This description is correct, but it is not yet very informative. If we stop at this description, a matrix looks like a table, and matrix algebra looks like a collection of rules for moving numbers around. In this chapter we will try to see more than that. A matrix is a way of organizing information. It can describe data, relations between quantities, coefficients of equations, or a rule that transforms one vector into another.

The first skill is therefore not calculation, but reading. When we see a matrix, we should ask what its shape is, what its entries represent, and how its rows and columns are meant to be interpreted. The same array of numbers may be only a table in one problem, but in another problem it may describe a transformation, a system of equations, or a geometric operation.

\begin{definitionbox}
A matrix with \(m\) rows and \(n\) columns is called an \(m\times n\) matrix. It has the form
\[
A=
\begin{pmatrix}
a_{11} & a_{12} & \cdots & a_{1n}\\
a_{21} & a_{22} & \cdots & a_{2n}\\
\vdots & \vdots & \ddots & \vdots\\
a_{m1} & a_{m2} & \cdots & a_{mn}
\end{pmatrix}.
\]
The entry \(a_{ij}\) is located in row \(i\) and column \(j\).
\end{definitionbox}

This notation is simple, but it contains an important discipline. The first index tells us where to look vertically, and the second index tells us where to look horizontally. A student who confuses these indices will often make mistakes later when multiplying matrices or describing systems of equations.

\begin{examplebox}
Consider the matrix
\[
M=
\begin{pmatrix}
4 & -2 & 0\\
7 & 1 & 5
\end{pmatrix}.
\]
This is a \(2\times 3\) matrix. The entry \(m_{11}\) is \(4\), the entry \(m_{12}\) is \(-2\), and the entry \(m_{23}\) is \(5\). We are not doing any operation yet. We are only learning to read the object.
\end{examplebox}

The size of a matrix is part of its identity. A row matrix, a column matrix, and a square matrix may contain similar numbers, but they play different roles. A column matrix can represent a vector. A square matrix can sometimes represent a transformation from a space to itself. A rectangular matrix may describe a relation between two collections of quantities of different sizes.

\begin{definitionbox}
A matrix is called square if it has the same number of rows and columns. An \(n\times n\) matrix is a square matrix of order \(n\).
\end{definitionbox}

We will often begin with square matrices because many basic ideas are easiest to see there. However, rectangular matrices are just as important. The point is not to memorize a list of shapes, but to develop the habit of asking: what kind of object am I holding, and what operations make sense for this object?
